\section{Knowing What You Do Not Know}

\subsection{The Map of Uncertainty}

First-time founders often believe the goal of this chapter is to construct an impressive personal narrative. The far more valuable output is a rigorous \textbf{map of uncertainty}.

A typical founder knows some things very well, knows some things indirectly, and has blind spots they cannot yet perceive. \textbf{Markets punish blind spots most severely because ignorance makes bad assumptions look obvious.}

I would maintain 4 separate registers:

\begin{enumerate}[label=\textbf{\arabic*.},leftmargin=11pt,itemsep=1.5pt,topsep=1.5pt,parsep=0pt]
  \item \textbf{Demonstrated assets:} Capabilities, relationships, and knowledge verified by E4 or E5 evidence.
  \item \textbf{Promising but unverified assets:} Hypotheses likely to help, but untested under real commercial pressure.
  \item \textbf{Known gaps:} Critical capabilities, buyer access, or balance sheet reserves you already know are absent.
  \item \textbf{Unknown unknowns to investigate:} Operational territories where you do not yet understand the landscape enough to formulate specific hypotheses.
\end{enumerate}

Consider an AI engineer proposing to train finance analysts:
\begin{itemize}[leftmargin=11pt,itemsep=1.5pt,topsep=1.5pt,parsep=0pt]
  \item \textit{Demonstrated:} Building and evaluating production machine-learning models.
  \item \textit{Unverified:} Explaining statistical architectures to finance teams without formal math backgrounds.
  \item \textit{Known gap:} Access to budget holders, compliance officers, and enterprise buyers at major banks.
  \item \textit{Unknown unknowns:} Statutory procurement cycles, regulatory certifications, data privacy constraints, and analyst shift demands.
\end{itemize}

The founder who makes these distinctions explicit can run a focused discovery process. The founder who conflates all 4 into ``my unfair advantage'' has set a dangerous trap.

\subsection{The Falsification Question}

For each central advantage, ask:

\begin{tcolorbox}[
  enhanced,
  colback=BayesPale!25!white,
  colframe=BayesBorder!60!white,
  boxrule=0.4pt,
  arc=2pt,
  left=6pt,right=6pt,top=4pt,bottom=4pt,
  before skip=3pt,after skip=4pt
]
\footnotesize\color{BayesInk}
\textbf{If this advantage were mostly imaginary, what concrete observation would reveal that within 2 weeks?}
\end{tcolorbox}

If you think you have buyer access, request 3 introductions. If you think your pedagogy is superior, run an independent learner assessment on unseen material. If you think you can build cheaply, track full end-to-end delivery hours. If you think a problem is urgent, examine what customers currently spend in cash or wasted labor to cope with it.

What you choose to measure reveals whether you genuinely want to learn or merely seek reassurance.
